\section{C. Integral presentations and their fppf meaning}
\label{sec:C}
The relevant question is C3, scans 334--336, manuscript folios
467--468. Scan 336 continues folio 468; its archive footer 335 is
not a manuscript folio. All coproducts and normal subgroup sheaves
below are on the big fppf site. Put

$$
s=\begin{pmatrix}0&-1\\1&0\end{pmatrix},\quad
 r=\begin{pmatrix}0&1\\-1&1\end{pmatrix},\quad
 u(t)=\begin{pmatrix}1&-t\\0&1\end{pmatrix},\quad
 v(t)=su(t)s^{-1}=\begin{pmatrix}1&0\\t&1\end{pmatrix}.
$$

Thus $s^2=r^3=-I$, $u(1)=sr$, and $r^{-1}su(1)=-I$.
Underlines denote constant group sheaves.

\begin{statement}{C.1}{Proved}{The actual source}
For $\Gamma_s=\operatorname{SL}_2(\mathbb Z)$ and
$\Gamma_p=\operatorname{PSL}_2(\mathbb Z)$, the sources are

$$
A_s=\underline\Gamma_s*_{\underline{\mathbb Z}}\mathbb G_a,
 \qquad A_p=\underline\Gamma_p*_{\underline{\mathbb Z}}\mathbb G_a.
$$

The maps from $\mathbb Z$ send $n$ to $u(1)^n$ and $n\cdot1$.
They map respectively to $\operatorname{SL}_2$ and
$\operatorname{PGL}_2$; the signed source maps to
$\operatorname{GL}_2$ through $\operatorname{SL}_2$.
The finite factors are constant $C_4,C_6$ with common $C_2$,
or projectively constant $C_2,C_3$. They are not $\mu_4,\mu_6$.
\end{statement}
\begin{proof}
The matrix identities give the cusp relations; C.2 identifies the
abstract group of finite generators. The assertion then follows
from the pushout universal property. The distinction is already
visible over $\mathbb Q$: constant $C_4$ has four rational
elements, whereas $\mu_4$ has two. At two, the matrix maps from
the constant factors identify some sections because $-I=I$.
Thus these maps need not be closed embeddings in every fibre.
Replacing the factors by schematic matrix closures changes the
problem.
\end{proof}
For a connected ring, distinguish the word groups

$$
W_s(R)=\Gamma_s*_{\mathbb Z}(R,+),\qquad
 W_p(R)=\Gamma_p*_{\mathbb Z}(R,+)
$$

from the sheaf sections $A_s(R),A_p(R)$.

\begin{statement}{C.2}{Known}{Integral amalgams and the level-two subgroup}

$$
\operatorname{PSL}_2(\mathbb Z)=C_2*C_3,\qquad
 \operatorname{SL}_2(\mathbb Z)=C_4*_{C_2}C_6.
$$

The kernel of
$\operatorname{PSL}_2(\mathbb Z)\to\operatorname{SL}_2(\mathbb F_2)$
is free on $u(2),v(2)$.
See \citet[Theorems 1.1,3.1,C.1]{C_ConradSL2Z}.
\end{statement}
\begin{proof}
Euclidean row reduction gives generation by $s,u(1)$.
On the real projective line, $s(x)=-1/x$ sends positive numbers
to negative numbers, and both nonidentity powers of
$r(x)=1/(1-x)$ send negative numbers to positive numbers.
A reduced word can be conjugated either into a single factor
or into a word beginning and ending in $C_3$: the two nonidentity
elements of $C_3$ permit a choice avoiding cancellation.
The latter word sends the negative domain into the positive
domain. This proves the projective free product.

In the signed amalgam, $z=s^2=r^3$ is central of order two and
the quotient by $z$ is the projective free product.
The kernel of the map to $\operatorname{SL}_2(\mathbb Z)$
therefore lies in $\langle z\rangle$, and $z$ maps to $-I\ne I$.

The projective level-two kernel has index six and intersects
conjugates of the finite factors trivially. It acts freely on
the amalgam tree; its quotient has six edges and
$6/2+6/3=5$ vertices, hence free rank two.
A matrix in $\Gamma(2)$ has odd diagonal and even off-diagonal
entries. Even elementary row operations implement the Euclidean
algorithm using nearest even quotients. Equality in the usual
remainder bound is excluded by opposite parities, so absolute
values strictly decrease. The final matrix is an even upper
elementary matrix, up to a central sign. Consequently
$\Gamma(2)=\langle u(2),v(2),-I\rangle$.
The two projective generators therefore form a free basis.
\end{proof}

\begin{statement}{C.3}{Proved}{Sheaf surjectivity and determinant}
The signed map to $\operatorname{SL}_2$ is Zariski-epimorphic;
the projective map to $\operatorname{PGL}_2$ is fppf-epimorphic.
The signed map to $\operatorname{GL}_2$ has image
$\operatorname{SL}_2$ and determinant cokernel $\mathbb G_m$.
\end{statement}
\begin{proof}
For a unit $t$ set

$$
w(t)=u(t)su(t^{-1})s^{-1}u(t),\qquad h(t)=w(t)s^{-1}.
$$

Multiplication gives

$$
w(t)=\begin{pmatrix}0&-t\\t^{-1}&0\end{pmatrix},\qquad
 h(t)=\operatorname{diag}(t,t^{-1}).
$$

For $g=\left(\begin{smallmatrix}a&b\\c&d\end{smallmatrix}\right)$
of determinant one, the words

$$
W_a(g)=v(c/a)h(a)u(-b/a),\qquad
 W_c(g)=sv(-a/c)h(c)u(-d/c)
$$

give $g$ where $a$, respectively $c$, is invertible.
These opens cover. A projective point locally lifts to a
general-linear matrix. If its determinant is $d$, adjoining
$t^2=d^{-1}$ is finite free and faithfully flat and permits
rescaling to determinant one, including in characteristic two.
Finally determinant is split by $t\mapsto\operatorname{diag}(t,1)$.
\end{proof}

\begin{statement}{C.4}{Refuted}{The bare kernel is noncentral}
Over every nonempty base, the signed and projective maps have
nontrivial noncentral kernels. The proposed relation does not
present the target and does not have a central $K_2$ kernel.
\end{statement}
\begin{proof}
In characteristic zero let $H_0=\Gamma$, $C_0=\mathbb Z$.
In characteristic $p$ let

$$
H_p=\Gamma/\langle\!\langle u(1)^p\rangle\!\rangle,\qquad
 C_p=\mathbb F_p,
$$

with the signed or projective choice of $\Gamma$.
The cusp embeds in $H_p$, since its matrix image has order $p$.
Over a field $K$ the word source is $H_p*_{C_p}(K,+)$.
Reduced words remain reduced under field extension and survive
fppf sheafification: a supposed local trivialization would,
on passing to a residue field of a nonempty covering member,
contradict this normal form.

Over $K=k(t,v)$ consider

$$
K(t)=w(t)u(1)w(t)^{-1}su(-t^{-2})s^{-1}.
$$

Its matrix is one, since $w(t)u(1)w(t)^{-1}=v(t^{-2})$.
Its reduced syllables are

$$
u(t),s,u(t^{-1}),s^{-1}u(1)s,u(-t^{-1}),s^{-1},
 u(-t),s,u(-t^{-2}),s^{-1}.
$$

Every additive syllable lies outside the prime cusp; every
constant syllable does too, as its matrix shows.
This is a cyclically reduced word of length ten.
Its commutator with $u(v)$ reduces only by combining the final
additive syllables to $u(-t-v)$, leaving length twenty-one.
Thus the sheaf kernel is noncentral. Every nonempty base
admits this test after passage to a residue field and a
transcendental extension.

Central-extension results with additional root relations do
not apply to this different source; compare
\citet[Introduction and Theorem 4.7]{C_BrylinskiDeligne2001}.
\end{proof}

\begin{statement}{C.5}{Proved}{The exact kernel by nine relation families}
Let $U_a,U_c$ be the two opens of C.3.
The signed kernel is the normal subgroup sheaf generated by

$$
W_a(g)W_c(g)^{-1}\quad(g\in U_a\cap U_c)
$$

and the eight families

$$
W_i(g)W_j(g')W_k(gg')^{-1},\qquad
 i,j,k\in\{a,c\},\quad g\in U_i,\ g'\in U_j,\ gg'\in U_k.
$$

Projectively, additionally kill $h(z)$ for every section of
the full schematic $\mu_2$. These statements commute with
every base change, including $\mathbb Z_2,\mathbb Z_3$.
\end{statement}
\begin{proof}
In the quotient $Q$ the overlap relation glues the words
to $\psi:\operatorname{SL}_2\to Q$.
The eight multiplication relations make $\psi$ a homomorphism,
since their domains cover all pairs.
The composite with the matrix map is the identity.
Conversely $W_a(u(t))=u(t)$ and $W_c(s)=s$;
$h(1)=1$ follows from the integral relations.
Since $r=s^{-1}u(1)$, these sections generate the source as
a sheaf, proving the other composite is the identity.
Thus the stated normal subgroup is the entire kernel.

Passing from $A_s$ to $A_p$ kills the constant sign.
After imposing the Gauss relations, killing all schematic
$h(z)$ gives $\operatorname{SL}_2/\mu_2=\operatorname{PGL}_2$.
At two the constant sign alone does not kill the infinitesimal
centre. The constructions and inverses are integral, proving
base-change compatibility.
\end{proof}

\begin{statement}{C.6}{Proved}{The general-linear presentation}
Adjoin a multiplicative generator $D(a)$, $a\in\mathbb G_m$,
to the signed quotient in C.5, with

$$
D(a)u(t)D(a)^{-1}=u(at),\qquad D(a)sD(a)^{-1}=w(a).
$$

The result is $\operatorname{GL}_2$ over every base.
\end{statement}
\begin{proof}
These are conjugation by $\operatorname{diag}(a,1)$ on
generators. The resulting semidirect product is identified
by $(g,a)\mapsto g\operatorname{diag}(a,1)$, with inverse

$$
M\longmapsto
 \bigl(M\operatorname{diag}(\det M,1)^{-1},\det M\bigr).
$$

\end{proof}

\begin{statement}{C.7}{Proved}{Orders of the word groups}

$$
\begin{array}{c|rrrrrrr}
R&\mathbb F_2&\mathbb F_3&\mathbb F_4&\mathbb F_5&\mathbb F_7&
\mathbb Z/4&\mathbb Z/9\\\hline
W_p(R)&6&12&\infty&60&\infty&24&\infty\\
W_s(R)&6&24&\infty&120&\infty&48&\infty
\end{array}
$$

Here infinity means countably infinite.
At $\mathbb Z$ the word groups are exactly those of C.2;
this does not identify them with all sheaf sections.
\end{statement}
\begin{proof}
For cyclic additive group of order $n$ the presentations are

$$
H_n^p=\langle s,r:s^2=r^3=(sr)^n=1\rangle,
$$


$$
H_n^s=\langle s,r:s^4=r^6=1,\ s^2=r^3,\ (sr)^n=1\rangle.
$$

The executed Todd--Coxeter program enumerates the trivial
subgroup for $n=2,3,4,5$, obtaining complete projective tables
of degrees $6,12,24,60$ and signed tables of degrees
$6,24,48,120$. It checks every relator at every coset.
New cosets are introduced only at undefined edges and
identifications arise only from relators and inverse-edge
consistency; completion gives upper bounds.
The independently enumerated matrix images over
$\mathbb F_2,\mathbb F_3,\mathbb Z/4,\mathbb F_5$
give matching lower bounds.

For $\mathbb F_4$ the source in both cases is
$\mathfrak S_3*_{C_2}(\mathbb F_4,+)$.
Both inclusions are proper, giving arbitrarily long reduced words.
For $n=7,9$, set

$$
t=2\cos(\pi/n),\quad d=\sqrt{t^2-3},\quad
 S=\begin{pmatrix}0&-1\\1&0\end{pmatrix},\quad
 R=\begin{pmatrix}1/2&(t+d)/2\\(d-t)/2&1/2\end{pmatrix}.
$$

Their traces give projective orders $2,3,n$ for $S,R,SR$.
Their upper-half-plane fixed points differ since $d>0$.
A finite subgroup has a common fixed point: average a positive
definite form over its finite lift to $\operatorname{SL}_2(\mathbb R)$.
Hence this image is infinite. Its signed cover is also infinite.
At $\mathbb Z$ the additive factor is the cusp itself.
No infinitude conclusion is based on an enumeration timeout.
\end{proof}

\begin{statement}{C.8}{Proved}{Actual sections over the finite rings}

$$
\begin{array}{c|rrrrrrr}
R&\mathbb F_2&\mathbb F_3&\mathbb F_4&\mathbb F_5&\mathbb F_7&
\mathbb Z/4&\mathbb Z/9\\\hline
A_p(R)&6&12&\aleph_0&\aleph_0&\aleph_0&\aleph_0&\aleph_0\\
A_s(R)&6&\aleph_0&\aleph_0&\aleph_0&\aleph_0&\aleph_0&\aleph_0
\end{array}
$$

Additional sections in the finite-word cases include
infinite-order descended matrix-kernel elements.
\end{statement}
\begin{proof}
For a finite field of characteristic $p$,

$$
A(\mathbb F_q)=
 \bigl(H_p*_{\mathbb F_p}(\overline{\mathbb F}_p,+)\bigr)
 ^{\operatorname{Frob}_q}.
$$

Indeed a section is locally a word; a closed point of an affine
fppf cover realizes it over a finite extension.
C.4 supplies the needed injections. A word uses finitely
many algebraic coefficients, so invariance is precisely finite
Galois descent.

If $C$ is malnormal in $H$, an automorphism acting trivially
on $H$ and with fixed subgroup $C$ in the additive factor
$V$ has fixed subgroup $H$ in $H*_C V$.
To prove this, compare a reduced word with its transform.
Normal-form equality allows only cusp insertions between
syllables. Across $h\notin C$, a nontrivial insertion would
belong to $C\cap hCh^{-1}$, impossible.
Thus additive syllables are fixed and lie in $C$, contradicting
reducedness; a lone additive syllable follows by cancellation.
At two, $H=\mathfrak S_3$ and $C=C_2$ is malnormal.
Projectively at three, $H=\mathfrak A_4$ and $C=C_3$
is malnormal. This proves the finite exceptions.

For signed $p=3,5$, choose $a^p-a=1$ and $z=-I$.
Then

$$
u(a)zu(-a)z^{-1}
$$

is Frobenius fixed since $z$ commutes with the cusp.
It is cyclically reduced of length four and has infinite
order, with identity matrix.
Projectively at five use
$h=[\operatorname{diag}(2,3)]$, which inverts the cusp.
Then $u(a)hu(a)h^{-1}$ is Frobenius fixed, cyclically
reduced, and has identity matrix.
At $\mathbb F_4,\mathbb F_7$ the infinite word subgroup suffices.

For $\mathbb Z/4$, use the finite etale Galois cover

$$
R'=(\mathbb Z/4)[x]/(x^2+x+1),\quad
 \tau(x)=-x-1,\quad a=2x,\quad \tau(a)=a+2.
$$

Signed, use the same word with $z=-I$.
Projectively, use

$$
h=\left[\begin{pmatrix}1&0\\2&1\end{pmatrix}\right],
 \qquad u(a)hu(-a)h^{-1}.
$$

The bridges commute with $u(2)$, so the words descend.
They are cyclically reduced of length four with identity
matrix. In the projective calculation, the cross products
containing $2\cdot2x$ vanish.
They survive sheafification: $R'$ is local Artinian, and
any nonzero flat algebra over it is faithfully flat.
A nonempty localization of a covering member is still
faithfully flat. Hence $R'$ embeds, and $a$ cannot become
an integer. The syllables remain reduced.
The same reasoning embeds the infinite word subgroup
over $\mathbb Z/9$ into the sheaf sections, since faithful
flatness preserves characteristic exactly nine and its cusp.

Over a finite base there are countably many finitely presented
affine algebras and words. Quasi-compactness allows finite
refinements of covers. Thus there are at most countably
many sections; the explicit subgroups give equality.
\end{proof}

\begin{statement}{C.9}{Known}{Amalgams after inverting a prime}
For $d_p=\operatorname{diag}(1,p)$,

$$
\operatorname{SL}_2(\mathbb Z[1/p])
 =\operatorname{SL}_2(\mathbb Z)*_{\Gamma_0(p)}
 d_p\operatorname{SL}_2(\mathbb Z)d_p^{-1},
$$

where $\Gamma_0(p)$ consists of matrices with $p\mid c$.
In coordinates on the second copy the edge inclusion is

$$
\begin{pmatrix}a&b\\c&d\end{pmatrix}
 \longmapsto\begin{pmatrix}a&pb\\c/p&d\end{pmatrix}.
$$

Both indices are $p+1$. This presents the target, not
$\Gamma_s*_{\mathbb Z}\mathbb Z[1/p]$.
\end{statement}
\begin{proof}
Use the tree of homothety classes of $\mathbb Z_p$-lattices
in $\mathbb Q_p^2$. Determinant one preserves index parity,
so edges are not inverted. The stabilizers of
$\mathbb Z_p^2$ and $d_p\mathbb Z_p^2$ are the stated
groups: integrality at $p$ removes all allowed denominators.
Their intersection is $\Gamma_0(p)$.
Neighbours are the $p+1$ lines in $\mathbb F_p^2$.
Integral reduction is onto $\operatorname{SL}_2(\mathbb F_p)$,
which is transitive on lines. Induction along paths gives
edge transitivity for the group generated by the two
stabilizers. Every element of the full group then differs
from an element of that generated subgroup by an edge
stabilizer, proving generation.
A reduced alternating word moves the edge along a
nonbacktracking path. Thus the only identifications are
those of the intersection. Compare
\citet[I.1.2, Theorem 1; II.1.4, Theorem 3 and Corollary 2]{C_Serre1980}.

Substituting $t=p$ in C.4 gives the reduced syllables

$$
u(p)s,\ u(1/p),\ s^{-1}u(1)s,\ u(-1/p),\
 s^{-1}u(-p)s,\ u(-1/p^2),\ s^{-1}.
$$

The four constant syllables are outside the cusp and the
other three parameters are not integers.
This length-seven matrix-identity word is nontrivial and,
by passage to a generic field, survives sheafification.
\end{proof}

\begin{statement}{C.10}{Proved}{Matrix images and every missed word image}

$$
\begin{array}{c|rrrr}
R&|\operatorname{SL}_2(R)|&|\operatorname{GL}_2(R)|&
|\operatorname{im}W_p(R)|&|\operatorname{PGL}_2(R)|\\\hline
\mathbb F_2&6&6&6&6\\
\mathbb F_3&24&48&12&24\\
\mathbb F_4&60&180&60&60\\
\mathbb F_5&120&480&60&120\\
\mathbb F_7&336&2016&168&336\\
\mathbb Z/4&48&96&24&48\\
\mathbb Z/9&648&3888&324&648
\end{array}
$$

The signed word image is all of $\operatorname{SL}_2(R)$.
The projective word image is exactly the square-determinant
subgroup. The program lists every omitted point by canonical
matrix code. Actual sheaf images are the same projective
subgroups over the listed finite fields and $\mathbb Z/9$.
Over $\mathbb Z$ the actual image is
$\operatorname{PSL}_2(\mathbb Z)$. The full sheaf image over
$\mathbb Z/4$ is excluded from this assertion.
\end{statement}
\begin{proof}
Over $\mathbb F_q$, choose a nonzero first column in $q^2-1$
ways; determinant-one second columns number $q$, while
independent second columns number $q^2-q$.
Successive special-linear kernels over $\mathbb Z/p^r$
have $p^3$ elements, by the trace-zero first-order condition.
Thus the order is $p^{3(r-1)}|\operatorname{SL}_2(\mathbb F_p)|$.
Determinant multiplies it by $\varphi(p^r)$ for the
general-linear group. Divide by scalar units for PGL,
and by scalar solutions $z^2=1$ for the special-linear
projective image. This proves the table.

Over these local rings the Gauss words give every
special-linear matrix. Their Picard groups vanish, so
projective matrices have general-linear representatives.
A representative rescales to determinant one exactly
when its determinant is a square unit.
The program encodes a matrix by $a+qb+q^2c+q^3d$ and
takes the least code of its scalar multiples.
Its missed-point lists are therefore complete without
duplicates. For the original map to GL, the omitted
points are exactly those of determinant other than one.

For a Frobenius-fixed reduced word over $\mathbb F_q$,
normal-form equality implies $a^q-a\in\mathbb F_p$
for each additive coefficient: representatives change
only by cusp elements at their ends.
Iteration $p$ times gives $a\in\mathbb F_{q^p}$.
Choose determinant-one representatives of the finitely
many constant syllables; the matrix word is defined over
that extension. For odd $p$, a nonsquare stays nonsquare
in this odd-degree extension, so no additional projective
class can descend. In characteristic two every finite-field
unit is a square.
Modulo nine, a unit is square exactly when its reduction
modulo three is square. Reduction excludes the remaining
classes. Determinant-minus-one integral matrices likewise
reduce to nonsquare projective classes modulo three.
All integral determinant-one matrices already occur by C.2.
This does not identify the integral sheaf point group
with its word subgroup.
\end{proof}

\begin{statement}{C.11}{Open}{The precise remaining lifting question}
It remains undecided here whether

$$
[\operatorname{diag}(3,1)]\in
 \operatorname{PGL}_2(\mathbb Z/4)
$$

lifts through $A_p(\mathbb Z/4)$.
The image has order $24$ or $48$; lifting this one class
decides which. This is a residual of this work, not a
literature-wide openness claim.
Also unproved are equality of integral sheaf sections
with their word subgroups and full descriptions of the
descended Artinian point groups.

The corrected presentation theorem C.3--C.6 is complete:
the constant-factor sources are sheaf epimorphisms,
their noncentral kernels are the specified normal
subgroup sheaves of Gauss relations, and adjoining
$\mathbb G_m$ gives GL. It holds over every base.
\end{statement}
\begin{proof}[Partial result, executed computation, and residual]
The square-determinant subgroup has index two.
Over $R'=(\mathbb Z/4)[x]/(x^2+x+1)$ enumerate all reduced
words $u(a)hu(b)$ with nonconstant $a,b$ and
$h\in H_4^p\setminus C_4$. There are $12\cdot20\cdot12=2880$.
Normal form shows that such a word is fixed by $\tau$
exactly when

$$
c=\tau(a)-a\in\mathbb Z/4,\quad
 d=\tau(b)-b\in\mathbb Z/4,\quad u(c)hu(d)=h.
$$

The executed search finds precisely $64$ fixed words.
Encoding $a_0+a_1x$ as $a_0+4a_1$, both additive codes
independently range over $8,9,10,11$.
Bridge codes are $97,103,105,111$ in C.10's convention;
every combination occurs. Their matrix images are exactly
those four codes, all with square determinant.

Indeed the coefficient of $x$ must be two, so both Galois
differences are two. The four bridges outside the cusp
commuting with $u(2)$ are $u(j)h$, with $h$ as in C.8.
Their matrices also commute with $u(2x)$, giving the
four stated targets.
Arbitrary word lengths and arbitrary fppf covers,
including ramified covers at two, are still untreated.
C.8 already proves the actual point group infinite,
so substituting the finite word group cannot decide
the remaining question.
\end{proof}

\paragraph{Computation provenance.}
The original Python programs were executed before the
workspace became inaccessible. Their original bytes are
not represented as recovered. The replacement JavaScript
and its actual output supply the finite cosets, finite
matrix counts, missed word-image points, and the
2880-word descent search. Infinitude and noncentrality
are proved above, not inferred from finite tests.
